\documentclass[10pt,a4paper,twocolumn]{article}
\usepackage[utf8]{inputenc}
\usepackage[spanish]{babel}
\usepackage{geometry}
\geometry{margin=2cm, columnsep=0.8cm}
\usepackage{hyperref}
\usepackage{tikz}
\usetikzlibrary{arrows.meta, positioning}
\usepackage{booktabs}

\title{Actividad 1 --- Construir un grafo Web}
\author{Angel Armando Carreño González}
\date{\today}

\begin{document}
\maketitle

\section{Tema y páginas seleccionadas}
Tema: \textbf{Sistemas de Recuperación de Información (SRI)}.
Se seleccionaron 9 páginas relacionadas con el tema:

\begin{enumerate}
  \item[P1] Wikipedia -- \textit{Information retrieval}
  \item[P2] Documentación de Elasticsearch
  \item[P3] Apache Lucene
  \item[P4] Artículo original de PageRank (Stanford)
  \item[P5] Google Search Central (guía de posicionamiento)
  \item[P6] TREC (Text REtrieval Conference, NIST)
  \item[P7] arXiv -- categoría Information Retrieval
  \item[P8] Scikit-learn -- documentación Precision-Recall
  \item[P9] Libro \textit{Introduction to Information Retrieval} (Manning et al., Stanford)
\end{enumerate}

\section{Enlaces registrados}
Se registraron 18 enlaces dirigidos (una arista $A \to B$ significa que la página $A$ enlaza a $B$):

\begin{center}
\begin{tabular}{cl}
\toprule
\# & Enlace \\
\midrule
1 & P1 $\to$ P4 \\
2 & P1 $\to$ P6 \\
3 & P1 $\to$ P9 \\
4 & P2 $\to$ P3 \\
5 & P2 $\to$ P5 \\
6 & P3 $\to$ P2 \\
7 & P4 $\to$ P1 \\
8 & P4 $\to$ P5 \\
9 & P5 $\to$ P2 \\
10 & P5 $\to$ P4 \\
11 & P6 $\to$ P1 \\
12 & P6 $\to$ P7 \\
13 & P6 $\to$ P9 \\
14 & P7 $\to$ P6 \\
15 & P7 $\to$ P9 \\
16 & P8 $\to$ P1 \\
17 & P8 $\to$ P2 \\
18 & P9 $\to$ P1 \\
\bottomrule
\end{tabular}
\end{center}

\section{Grafo dirigido}
\begin{center}
\begin{tikzpicture}[node distance=1.4cm,
  nodo/.style={circle, draw, minimum size=0.85cm, font=\footnotesize},
  flecha/.style={-{Stealth}, thick}]
  \node[nodo] (P1) {P1};
  \node[nodo] (P2) [right=of P1] {P2};
  \node[nodo] (P3) [right=of P2] {P3};
  \node[nodo] (P4) [below=of P1] {P4};
  \node[nodo] (P5) [below=of P2] {P5};
  \node[nodo] (P6) [below=of P3] {P6};
  \node[nodo] (P7) [below=of P6] {P7};
  \node[nodo] (P8) [left=of P4] {P8};
  \node[nodo] (P9) [below=of P4] {P9};
  \draw[flecha] (P1) -> (P4);
  \draw[flecha] (P1) -> (P6);
  \draw[flecha] (P1) to[bend left=20] -> (P9);
  \draw[flecha] (P2) -> (P3);
  \draw[flecha] (P2) -> (P5);
  \draw[flecha] (P3) to[bend left=20] -> (P2);
  \draw[flecha] (P4) -> (P1);
  \draw[flecha] (P4) -> (P5);
  \draw[flecha] (P5) -> (P2);
  \draw[flecha] (P5) to[bend left=15] -> (P4);
  \draw[flecha] (P6) -> (P1);
  \draw[flecha] (P6) -> (P7);
  \draw[flecha] (P6) to[bend left=15] -> (P9);
  \draw[flecha] (P7) to[bend left=15] -> (P6);
  \draw[flecha] (P7) -> (P9);
  \draw[flecha] (P8) -> (P1);
  \draw[flecha] (P8) to[bend left=15] -> (P2);
  \draw[flecha] (P9) to[bend left=20] -> (P1);
\end{tikzpicture}
\end{center}

\section{Nodos con más enlaces}
\begin{center}
\footnotesize
\begin{tabular}{ccc}
\toprule
Nodo & Entrantes & Salientes \\
\midrule
P1 & 4 & 3 \\
P2 & 3 & 2 \\
P3 & 1 & 1 \\
P4 & 2 & 2 \\
P5 & 2 & 2 \\
P6 & 2 & 3 \\
P7 & 1 & 2 \\
P8 & 0 & 2 \\
P9 & 3 & 1 \\
\bottomrule
\end{tabular}
\end{center}

\textbf{Autoridades} (muchos entrantes): P1 (4), P2 y P9 (3). \\
\textbf{Hubs} (muchos salientes): P1 y P6 (3).

\section{Discusión: ¿muchos entrantes = la mejor página?}
No necesariamente. El conteo de entrantes es solo una señal:

\begin{itemize}
  \item \textbf{Calidad vs.\ cantidad:} un enlace desde P6 (TREC, fuente experta) vale más que varios desde páginas periféricas. Esta es la intuición de PageRank y HITS: importan quiénes enlazan, no solo cuántos.
  \item \textbf{Sesgo y spam:} una página puede acumular entrantes por navegación, directorios o granjas de enlaces, sin ser la más precisa ni completa.
  \item \textbf{Rol distinto:} P9 (libro) tiene 3 entrantes y 1 saliente: es autoridad temática, pero P6 es mejor hub para descubrir recursos. ``Mejor'' depende de la necesidad: aprender (autoridad) vs.\ explorar (hub).
  \item \textbf{Nodo aislado:} P8 no recibe enlaces pero aporta 2 salientes útiles; el indegree cero no significa inútil.
\end{itemize}

Conclusión: el indegree orienta, pero la relevancia real exige ponderar la calidad de las fuentes (PageRank), el contenido y la necesidad del usuario, igual que en evaluación SRI se combinan Precisión, Recall y ranking.

\end{document}
